\chapter{Discusi\'on}
\label{cap:discusion}

\section{Cobertura y calidad de c\'odigo}

Elevar el umbral autom\'atico de JaCoCo de 60\,\% a 70\,\% y mantener LINE
en 91{,}80\,\% (muy por encima del umbral) sugiere que la suite de pruebas
creci\'o en proporci\'on al c\'odigo nuevo de la Unidad~IV, no solo en
n\'umero absoluto (109~$\to$~205 pruebas): el crecimiento de BRANCH
(76{,}87\,\%~$\to$~79{,}39\,\%) fue menor que el de LINE, lo que es
consistente con la naturaleza t\'ipica de c\'odigo de validaci\'on
condicional (por ejemplo, en \texttt{sp\_registrar\_consulta\_validada} y
sus equivalentes en Java) que a\~nade ramas dif\'iciles de cubrir en su
totalidad. La cobertura sigue midiendo ejecuci\'on, no correcci\'on: el
hecho de que 0 hallazgos \texttt{SQL\_*} coexistan con 205 pruebas en verde
es evidencia complementaria, no una prueba formal de ausencia de defectos.

\section{Rendimiento}

Las diez corridas de k6 mantienen el mismo patr\'on cualitativo observado
en la Tercera Entrega (latencia mediana de un solo d\'igito de
milisegundos en condiciones de cach\'e caliente, throughput cercano a
92~req/s), con tasa de error 0{,}0\,\% sostenida al duplicar el n\'umero de
corridas (3+3~$\to$~5+5). El an\'alisis de Wilcoxon con tama\~no de efecto
de Rosenthal muestra que el efecto de la cach\'e no es uniforme entre
corridas ($r$ desde $-0{,}07$ hasta $-0{,}84$), lo que sugiere que factores
del entorno de ejecuci\'on (por ejemplo, el estado de la JVM o del pool de
conexiones al inicio de cada corrida) influyen en la magnitud observable
del beneficio de cach\'e, sin que eso invalide la conclusi\'on cualitativa
de que la cach\'e reduce la latencia. Frente a Putri, Hadi \& Ramdani
(2017)~\cite{putri2017perftesting} (cap\'itulo~\ref{cap:trabajos-relacionados}),
BIOPET reporta el mismo tipo de evidencia de \emph{benchmarking} pero
complementada con otras cinco dimensiones de calidad, no en aislamiento.

\section{Seguridad}

El resultado combinado de las tres t\'ecnicas de auditor\'ia (manual,
din\'amica, est\'atica) es consistente entre s\'i: ning\'una detect\'o
patrones de inyecci\'on SQL explotable, y el \'unico hallazgo de severidad
alta (CSRF deshabilitado) tiene una mitigaci\'on arquitect\'onica expl\'icita
y documentada (secci\'on~\ref{sec:csrf}), no ignorada. Comparado con el
trabajo de Yang et al. sobre vulnerabilidades JWT~\cite{yang2026tokentimebomb},
BIOPET no reclama haber verificado la ausencia de las categor\'ias de
vulnerabilidad que ese estudio descubre mediante \emph{fuzzing} dirigido a
la librer\'ia \texttt{jjwt}: esa clase de an\'alisis queda expl\'icitamente
fuera del alcance de las tres t\'ecnicas usadas aqu\'i.

\section{Usabilidad}

Con $n=18$ y media 74{,}44/100, BIOPET se ubica sobre el umbral de
referencia de 68 puntos de Bangor, Kortum \& Miller~\cite{bangor2008sus},
en la categor\'ia ``Bueno''. El intervalo de confianza se estrech\'o
respecto a la muestra de $n=10$ de la Tercera Entrega, pero como se declara
en el cap\'itulo~\ref{cap:pruebas-calidad}, ampliar la muestra no resuelve
por s\'i solo el sesgo de reclutamiento por conveniencia. El participante
de menor puntaje (sin experiencia web previa) refuerza, de forma
consistente con la literatura de usabilidad, la relevancia de una fase de
orientaci\'on breve para usuarios de perfil similar en trabajo futuro.

\section{Calidad}

La combinaci\'on de cobertura alta, cero hallazgos SQL y auditor\'ia de
seguridad multim\'etodo en verde sugiere un nivel de calidad de c\'odigo
consistente con las buenas pr\'acticas esperadas en un proyecto de este
tipo, enmarcado bajo ISO/IEC~25010 (secci\'on~\ref{sec:iso25010}). El
hallazgo de SEO~82/100 (por debajo del umbral 90) detectado en la primera
corrida de Lighthouse (2026-08-01) se document\'o como \'area de mejora
concreta en vez de minimizarse, y qued\'o corregido en la corrida definitiva
del 2026-08-18 (SEO~100/100 en las 12 corridas, ambos perfiles): a la fecha
de cierre de este informe, las cuatro categor\'ias de Lighthouse cumplen su
umbral.

\section{Despliegue y reproducibilidad}

El despliegue real en Render, con \emph{healthcheck} en estado \texttt{UP}
y HTTPS activo, junto con la imagen inmutable en GHCR y el archivado
permanente en Zenodo, cierra la brecha entre ``el sistema funciona en mi
m\'aquina'' y ``el sistema es reproducible y accesible por terceros'' que
persist\'ia en la Tercera Entrega. La reproducibilidad de punta a punta
(\texttt{make up}, im\'agenes fijadas por \emph{digest}, migraciones desde
base vac\'ia verificadas) es, junto con la amplitud de evidencia
multim\'etrica, uno de los dos pilares que distinguen a BIOPET frente a
los trabajos relacionados revisados (cap\'itulo~\ref{cap:trabajos-relacionados}),
ninguno de los cuales documenta reproducibilidad de infraestructura con el
mismo nivel de detalle.

\section{Comparaci\'on con trabajos relacionados: l\'imites de la afirmaci\'on}

Ninguna comparaci\'on de este cap\'itulo implica superioridad absoluta de
BIOPET sobre Rodr\'iguez et al.~\cite{rodriguez2024oscrum}, Cede\~no Ochoa
et al.~\cite{cedeno2021veterinaria}, Putri et al.~\cite{putri2017perftesting}
o Yang et al.~\cite{yang2026tokentimebomb}: no existe una comparaci\'on
emp\'irica directa (mismos datos, mismo entorno) entre BIOPET y ninguno de
esos sistemas. La \'unica afirmaci\'on sostenible, con la evidencia
disponible, es que BIOPET reporta un conjunto \textbf{m\'as amplio} de
dimensiones de calidad medidas sobre un mismo artefacto que cualquiera de
los trabajos revisados individualmente, no que su calidad relativa en cada
dimensi\'on individual sea superior.

\section{Respuesta expl\'icita a las preguntas de investigaci\'on}
\label{sec:respuestas-rq}

\textbf{RQ1 --- \textquestiondown En qu\'e medida BIOPET satisface los
requisitos definidos para la Entrega Final?} El n\'ucleo \emph{Must Have}
(autenticaci\'on, autorizaci\'on por rol/propiedad, CRUD de mascotas) est\'a
completamente verificado, y los tres requisitos formalizados en la
Unidad~IV (\texttt{REQ-F-023} a \texttt{REQ-F-025}: gesti\'on
administrativa de usuarios, vacunas e integraci\'on de API externa) est\'an
verificados o implementados seg\'un el caso (cap\'itulo~\ref{cap:requisitos}). Un
subconjunto expl\'icito de requisitos \emph{Should}/\emph{Could} heredados
de la Entrega~1A (historial cl\'inico longitudinal completo, calendario
interactivo de citas, IoT, facturaci\'on digital, reportes exportables,
recomendaciones por IA) permanece \textbf{pendiente}, sin que este informe
lo declare cumplido. Respuesta: BIOPET satisface el n\'ucleo de requisitos
comprometido para esta entrega y una parte de los requisitos ampliados de
Unidad~IV; no satisface la totalidad del alcance funcional original de la
Entrega~1A.

\textbf{RQ2 --- \textquestiondown Qu\'e resultados presenta BIOPET en las
evaluaciones emp\'iricas ejecutadas?} 205 pruebas automatizadas (0
fallos/errores); cobertura JaCoCo 91{,}80\,\%~LINE/79{,}39\,\%~BRANCH sobre
un umbral autom\'atico de 70\,\%; diez corridas de rendimiento k6 con
0{,}0\,\% de tasa de error; usabilidad SUS con $n=18$, media 74{,}44/100
(categor\'ia ``Bueno''); calidad web Lighthouse con las cuatro categor\'ias
cumplidas en la corrida del 2026-08-18 (SEO~100/100); y seguridad auditada
por tres m\'etodos independientes (OWASP manual, ZAP din\'amico, SpotBugs
est\'atico), con 0 alertas de severidad alta y 0 hallazgos de inyecci\'on
SQL (cap\'itulos~\ref{cap:pruebas-calidad} y \ref{cap:seguridad}). Respuesta:
en las seis dimensiones evaluadas, BIOPET cumple los umbrales configurados
por el propio equipo, con las limitaciones de alcance documentadas en el
cap\'itulo~\ref{cap:amenazas-limitaciones-final} (entorno acad\'emico,
tama\~no muestral, mecanismo de sesi\'on de la corrida Lighthouse del
2026-08-18 sin documentar en prosa).

\textbf{RQ3 --- \textquestiondown El despliegue, CI, GHCR y archivado
permiten reproducir y verificar el artefacto?} S\'i, con evidencia real:
seis \emph{jobs} de CI en verde, despliegue activo en Render con
\emph{healthcheck} \texttt{UP}, imagen de contenedor publicada en GHCR con
referencia inmutable por \emph{digest}, y archivado permanente en Zenodo
con DOI real para el software y para el conjunto de datos de evidencias
(cap\'itulo~\ref{cap:despliegue}). La reproducibilidad tiene l\'imites
declarados: el tag Git \texttt{v1.0.0} y el GitHub Release a\'un no existen,
y el plan gratuito de PostgreSQL en Render no permite verificar hoy una
operaci\'on sostenida m\'as all\'a de 30 d\'ias
(secci\'on~\ref{sec:limitacion-30-dias}). Respuesta: el artefacto es
reproducible de punta a punta en el entorno de desarrollo y verificable en
producci\'on dentro de esas condiciones expl\'icitas.
